\section{Introduction}
The human sacroiliac joint (SIJ) represents an important mechanical compromise in hominin evolution: it must provide stable transmission of ground reaction forces during bipedal locomotion while simultaneously accommodating the compliance required for cephalopelvic clearance during parturition \citep{Washburn1960,Ruff2017,Pavlicev2020,Whitcome2017}. Biomechanically, this dual demand is governed by the principles of \emph{form closure} and \emph{force closure} \citep{Vleeming2012,Snijders1993}. Form closure is conferred by the interlocking wedge-like architecture of the sacrum within the pelvic ring, characterized by biconcave--biconvex articular congruency, macroscopic chondral ridges, and high frictional resistance \citep{Vleeming2012,Kiapour2019}. Force closure is provided dynamically by surrounding musculature and passively by an extensive ligamentous complex, including the anterior, interosseous, and dorsal sacroiliac ligaments, as well as the sacrotuberous and sacrospinous ligaments \citep{Snijders1993,Hammer2013,Hammer2019}.

Under upright axial gravity loading, the functional stability of the pelvic ring relies on \emph{nutation}---the anterior-inferior rotation of the sacral promontory relative to the ilia \citep{Vleeming2012,Snijders1993,Keizer2019}. Nutation increases tension across the primary posterior ligamentous stabilizers, drawing the auricular surfaces into tight apposition and self-bracing the articulation against vertical shear \citep{Snijders1993,Vleeming2012}. During cyclical gait, unilateral single-leg support (the single-stance phase) amplifies this demand: ground reaction forces are transmitted through a single acetabulum while body weight acts eccentrically on the contralateral side, generating pelvic tilt and asymmetric shear across the pelvic arch \citep{Bergmann2001JBiomech,Kutzner2016PLOS,Hammer2013}.

In contrast to upright locomotor support, parturition imposes an opposing mechanical demand: circumferential and transverse expansion of the pelvic cavity \citep{Pavlicev2020,BiomechPregnancy2022,AJOG2022Uterine}. As the fetal head traverses the inlet, midplane, and outlet, the surrounding pelvic ring is subjected to internal outward pressures and soft-tissue tensions that challenge sacroiliac joint congruency \citep{BiomechPregnancy2022,Fuchs2013PLOS,Almeida2024Uterus}. While pregnancy-related hormones such as relaxin contribute to systemic ligamentous laxity in vivo, the underlying bony and articular geometry of the pelvic ring fundamentally determines the passive kinematic pathways through which external expansion loads are absorbed \citep{Ruff2017,Betti2020,Weber2019,Weiss2020}. Although in vivo experimental studies consistently describe the SIJ as a low-motion joint (sub-millimetric translations and rotations rarely exceeding a few degrees; \citep{Goode2008,Meijer2014,Keizer2019,Noe2020}), they cannot isolate how differing loading directions selectively re-route 3D joint displacement vectors across diverse pelvic morphologies.

Finite-element (FE) analysis provides a valuable methodology to resolve these joint-level kinematics under standardized boundary conditions \citep{Anderson2020,Yang2020,Song2021,Calixte2022}. However, two fundamental questions remain unresolved in pelvic functional anatomy. First, how are standardized 3D force vectors partitioned into local joint degrees of freedom? Specifically, does simulated obstetric loading simply attenuate overall micromotion, or does it execute a qualitative \emph{directional re-routing} from sagittal nutation into transverse ring expansion? Second, what is the primary determinant of between-subject kinematic variability in the human population? Population variance in joint micromotion could plausibly originate from two distinct mechanistic ``channels'': macro-architectural geometric variation (pelvic proportions, inlet breadth, subpubic angle, and articular orientation; \citep{Ruff2017,Betti2020,Weber2019}) or tissue-level material heterogeneity (heterogeneous bone mineralization and apparent density; \citep{CarterHayes1977,Keller1994,Morgan2003Trabecular,Frost1987,Huiskes1987}). Establishing whether anatomy or bone tissue quality governs SIJ kinematics is important for evolutionary morphology, biomechanical assessment of peripartum pelvic girdle pain, and orthopedic surgical fixations \citep{Hammer2019,Kiapour2019}.

To resolve these questions, we investigated a standardized finite-element repository comprising 278 human pelves (150 female, 128 male; age 16--91 years) parameterized on a common reference template via an implicit continuous displacement mapping \citep{HenyssKuchar2025BoneDat}. The dataset features fully coupled, shape-only, and material-only simulation variants evaluated across five standardized static loading regimes spanning bipedal locomotion and multi-axial parturition proxies. We tested four anatomically grounded hypotheses:

\subsection*{Hypotheses}
\begin{enumerate}
\item \textbf{H1 (Load-Dependence and Stance Asymmetry):} Bilateral and unilateral upright stance will engage physiological self-bracing nutation ($\theta_{\text{nut}}$), with unilateral stance generating substantial bilateral kinematic asymmetry due to asymmetric acetabular support. Conversely, parturition-motivated mechanical proxies will de-emphasize sagittal nutation and vertical shear in favor of transverse compliance.
\item \textbf{H2 (Directional Contrasts in Kinematic State Space):} Parturition-motivated configurations will recruit directional displacements distinct from locomotor stance, shifting the dominant motion from craniocaudal shear ($\Delta\text{CC}$) to transverse outward distraction ($\Delta\text{ML}$) and sagittal translation.
\item \textbf{H3 (Sex-Related Differences and Morphometric Associations):} Female pelves will exhibit greater translational compliance under parturition-motivated configurations after adjusting for age and pelvic size, with this difference being attenuated when conditioning on transverse pelvic architecture (particularly subpubic arch angle and biischiadic width).
\item \textbf{H4 (Relative Contributions of Shape versus Material Heterogeneity):} Between-subject kinematic variance across the population will be governed primarily by macroscopic skeletal geometry rather than bone mineral density heterogeneity.
\end{enumerate}
As an exploratory secondary objective, we also evaluated allometric scaling relationships between pelvic volume and SIJ micromotion across all load configurations.
